% Round 18, root branch. Independently reviewed; see the bound audit.
\section{Smoothings of products of nonlinear hypersurfaces}
\label{r18nf:section}

The divisor-generation argument for the resolved hyperplane model
extends to smooth hypersurface factors of arbitrary positive degree.
It is useful that this proof does not require middle positivity on
the components: in dimension four the original components can
themselves be high-degree fourfold hypersurfaces.

\subsection{Input and the model}

Fix a finite valued extension $K/\mathbf Q_p$ with valuation ring
$R$, uniformizer $\pi$ and finite residue field $k$.  Let $n\ge4$,
let $m\ge3$, and choose forms
\[
 f_i\in R[X_0,\ldots,X_{n+1}]_{d_i},\qquad d_i\ge1,
 \quad 1\le i\le m,
 \qquad D=\sum_{i=1}^m d_i.
 \tag{NF18.1}\label{r18nf:factors}
\]
Assume that the hypersurfaces $H_i=V(\bar f_i)$ form a simple
normal-crossings divisor in $\mathbf P^{n+1}_k$: every intersection
$H_I$ with $1\le |I|\le n+1$ is a smooth complete intersection of
the expected dimension, and every intersection of $n+2$ factors is
empty.  In particular all positive-dimensional $H_I$ are geometrically
integral; zero-dimensional intersections are finite etale schemes
and need not consist of rational points.

Choose a degree-$D$ form $Q$ over $R$ whose reduction cuts every
positive-dimensional $H_I$ transversely and avoids every
zero-dimensional $H_I$.  Assume that
\[
              X=V(f_1\cdots f_m+\pi Q)\subset\mathbf P^{n+1}_K
 \tag{NF18.2}\label{r18nf:generic}
\]
is smooth and geometrically integral.  All conditions are geometric
conditions on the displayed actual reductions and on this generic
fiber.  No unspecified alteration or auxiliary extension is part
of the hypothesis.

Blow up the ambient $\mathbf P^{n+1}_R$ successively along the
strict transforms of the special divisors $H_1,\ldots,H_{m-1}$,
and let $\mathcal X$ be the strict transform of the hypersurface.
This is the identical ordered construction used in HP18, with
smooth equations $f_i$ in place of linear equations.

\begin{lemma}[Strata and connected centers]
\label{r18nf:model}
The resulting $\mathcal X$ is a regular integral projective
strictly semistable model of $X$, with reduced fiber
$\bigcup_{i=1}^m D_i$ and no additional components.
For every nonempty $I$ and $h=\max I$, its closed stratum $D_I$
is obtained from $H_I$ by the reverse-order blowups, for
$j=m,m-1,\ldots,h+1$, of the strict transforms of
\[
                    Z_{I,j}=H_I\cap H_j\cap V(\bar Q).
 \tag{NF18.3}\label{r18nf:centers}
\]
Empty centers are omitted.  Each occurring center is smooth,
and each of its successive strict transforms is isomorphic to
the original center.  In particular, for an edge $i<j$, every
center $Z_{ij,k}$ with $k>j$ has dimension $n-3\ge1$ and is
geometrically integral.
\end{lemma}

\begin{proof}
At a point of the reduced fiber the factor equations passing
through it are part of an etale coordinate system, by the
normal-crossings hypothesis.  If $\bar Q$ vanishes there,
transversality to their common stratum supplies one further
coordinate $y=-Q/u$, with $u$ the product of the unit factors.
The actual local equation is $x_1\cdots x_s=\pi y$.
If $\bar Q$ is nonzero, the equation is already a regular
strict crossing.  HP18's saturated ordered charts therefore
apply word for word; their derivation uses the independent
local parameters, not the linearity of the global $f_i$.
They prove regularity, reducedness, smoothness of all closed
strata and absence of new components.  Integrality, flatness
and projectivity follow as in that calculation.

On $H_I$ the reverse-order chart ideals are
$(\bar Q,\bar f_j)$.  At an overlap, changing trivializations
of their line bundles multiplies these generators by units,
so the local identifications glue to the global blowups in
\eqref{r18nf:centers}.  Each center is a smooth complete
intersection of dimension $n-|I|-1$ by hypothesis.
A previously blown center induces on a later center the
Cartier equation of the additional $H_j$.  That equation is
a nonzerodivisor whenever the intersection is nonempty,
because the positive-dimensional complete intersections are
geometrically integral and have the stated dimensions.
The strict transform is its blowup in a Cartier ideal,
hence is unchanged as a variety.  The local coordinate
calculation also proves it is still a regular codimension-two
center in the successive stratum.  Zero-dimensional centers
are disjoint from the other centers, since $Q$ avoids all
zero-dimensional $H_I$.

For an edge the center is a smooth positive-dimensional
complete intersection in projective space.  Such a complete
intersection is geometrically connected, and smoothness then
makes it geometrically integral.  This connectedness matters:
it gives exactly one exceptional divisor class per triple,
rather than an unaccounted collection of classes.
\end{proof}

\subsection{The native divisor matrix}

Let $h_{ij}$ be the pullback of the hyperplane class to
$E_{ij}=D_i\cap D_j$, $i<j$.  Let $e_{ij,k}$ be the total
transform of the exceptional divisor of the blowup along
$Z_{ij,k}$, for $k>j$.  The later centers are not contained
in this exceptional divisor, so this is also its strict
transform class.  Weak Lefschetz gives
$H^2((H_i\cap H_j)_{\bar k},\mathbf Q_\ell(1))
=\mathbf Q_\ell h$ because $\dim(H_i\cap H_j)=n-1\ge3$.
For any native divisor class, its geometric cycle class is a scalar
multiple of $h$.  Pairing with $h^{n-2}$ shows that this scalar
is rational.  Subtract that rational multiple of the actual
hyperplane divisor; the difference is homologically zero and
therefore numerically zero.  The ample class is numerically
nonzero.  Consequently the native numerical divisor group
is exactly $\mathbf Qh$.  The blowup divisor formula now gives
\[
 N^1(E_{ij})_{\mathbf Q}
   =\mathbf Q h_{ij}\oplus
                 \bigoplus_{k>j}\mathbf Q e_{ij,k}.
 \tag{NF18.4}\label{r18nf:picard}
\]
This equality is native: the centers are geometrically integral
and defined over $k$, and all listed divisor classes are native.

\begin{lemma}[The triple relations remove every exceptional divisor]
\label{r18nf:triple}
For $i<j<k$, DG18's actual triple relation has coordinates
\[
 \tau_{ijk}=
 \bigl(d_k h_{ij}-e_{ij,k},\ -d_j h_{ik},\ d_i h_{jk}\bigr)
 \tag{NF18.5}\label{r18nf:triple-formula}
\]
on its three incident edges, and zero elsewhere.  The images
of all $h_{ij}$ form a basis of
\[
              W^1/\operatorname{im}\tau.
 \tag{NF18.6}\label{r18nf:quotient-basis}
\]
In particular this quotient has dimension $\binom m2$.
\end{lemma}

\begin{proof}
On $H_i\cap H_j$, the triple divisor is $H_k$ and has
class $d_kh$.  The center $Z_{ij,k}$ is contained in this
divisor with multiplicity one.  Its blowup therefore changes
the strict-transform divisor class to $d_kh-e_{ij,k}$.
For the earlier steps, indexed by $l>k$, the center meets
that divisor properly and is not contained in it.  For a later
step $j<l<k$, the local ideals before the $k$-step are
$(x_k,y)$ and $(x_l,y)$.  In the blowup of $(x_k,y)$, the
strict divisor $x_k=0$ lies in the $y$-chart, with equation
$x_k/y=0$; the strict later center lies in the $x_k$-chart,
with equations $(x_l,y/x_k)=0$.  These loci are disjoint,
since the two displayed ratios are inverse on the overlap.
They remain disjoint under the subsequent center blowups.
Thus no other step subtracts an additional exceptional
divisor term.  These are identities of actual Cartier divisor
classes, with multiplicities checked by the independent
local parameters $f_k,Q$.

On $H_i\cap H_k$, no center indexed by $j<k$ is blown up.
All centers actually blown up meet $H_j$ properly, so its
strict-transform divisor class is just $d_jh_{ik}$.
The same reasoning on $H_j\cap H_k$ gives $d_ih_{jk}$.
The orientation signs are $+,-,+$, as fixed in DG18.
This proves \eqref{r18nf:triple-formula}.

In the direct sum of the bases
\eqref{r18nf:picard}, each triple column has exceptional
coefficient $-1$ at its unique $e_{ij,k}$ and zero at
every other exceptional basis element.  The exceptional
block is thus $-I_{\binom m3}$.  Eliminating these columns
removes all exceptional generators and introduces no
relation among the $h_{ij}$.  This proves the claimed basis.
The cancellation $\sigma_2\tau=0$ is the actual triple
Gysin cancellation of DG18 and does not discard higher strata.
\end{proof}

Let $\mathcal L$ be any relatively ample line bundle on
$\mathcal X$, and let $h$ denote the pullback of the ambient
$\mathcal O(1)$.  For every sufficiently large integer $M$,
the rational divisor class
\[
                         H_M=c_1(\mathcal L)+Mh
 \tag{NF18.7}\label{r18nf:ample}
\]
is relatively ample.  One may use a relatively ample bundle
from the projective blowup construction; no independently
chosen polarizations on different components are needed.
In the basis \eqref{r18nf:quotient-basis}, the classes given
by $H_M$ supported on each separate edge are the columns
of $M I+C$, for a fixed rational matrix $C$.
Its determinant is monic of degree $\binom m2$, and hence
is nonzero for all sufficiently large integers $M$.
It follows that, for such one global ample $H_M$,
\[
                       W^1=H_M W^0+\operatorname{im}\tau.
 \tag{NF18.8}\label{r18nf:generation}
\]

\subsection{Local Green consequences}

\begin{theorem}[Nonlinear-factor local exactness and strong G]
\label{r18nf:theorem}
For the model in Lemma~\ref{r18nf:model},
\[
                       U_2=0,
 \qquad
 K_{\mathrm{comp},\ell}^{\,n-2}=\ker\rho_{n-2}.
 \tag{NF18.9}\label{r18nf:vanishing}
\]
Consequently every specified geometrically homologically
trivial rational codimension-two cycle on $X$ satisfies
strong G in the precise ordinary-actual sense of the question.
If $n=4$, every geometrically homologically trivial rational
cycle satisfies G in every codimension.
The assertions reapply to all new cycles after each
unramified finite valued extension.
\end{theorem}

\begin{proof}
Equation \eqref{r18nf:generation} and the actual inclusion
$\operatorname{im}\tau\subseteq\ker\sigma_2$ verify the
divisor-generation hypothesis of DG18.  Its proof gives
$U_2=0$ using only divisor Hodge index.  In particular it
uses no middle numerical Hodge-standard hypothesis on
the possibly high-degree components $H_i$ or $D_i$.

We check the separate representative-lifting gate.  On the
smooth complete-intersection threefold or higher-dimensional
variety $H_i\cap H_j$ of dimension $s=n-1\ge3$, weak
Lefschetz gives $H^2=\mathbf Q_\ell h$.  Poincare duality
then makes $H^{2s-2}$ one-dimensional; the nonzero degree of
$h^s$ identifies its generator, so
$H^{2s-2}(s-1)=\mathbf Q_\ell h^{s-1}$.
A rational one-cycle has zero geometric class if and only if
its degree against $h$ is zero, because the top-degree
trace of $h^s$ is nonzero.  Numerical triviality therefore
implies zero geometric cycle class.  CD17's actual
curve-degree blowup decomposition preserves this implication
under every specified smooth blowup, including centers of
dimension greater than two: the only new curve-degree
coefficient is a zero-cycle on the center, whose class is
its degree.  Thus every resolved double $E_{ij}$ has the
required curve comparison.  CL17 proves the second
equality in \eqref{r18nf:vanishing}.

DQ17 now makes the refined quotient in adjoint degree
$a=n-2$ zero.  ER17 and VT17 give obstruction zero, and
the frozen reconstruction produces the form for the
original specified cycle on its exact support complement
against every compact ordinary actual test.  The ambient
resolution is not used as an analytic pushforward.

For $n=4$, codimension two was the only missing degree
beyond the frozen divisor, one-cycle and zero-cycle
theorems.  Equivalently, on every threefold double all
degrees needed by CL17 have the just-proved curve comparison
or the standard divisor and degree comparisons.  The
complete local conclusion follows.  Under an unramified
extension, the same equations, smooth strata and model
calculations apply after residue extension.  Reapply the
theorem to each cycle over the new field; there is no
assumption that it was pulled back from $K$.
\end{proof}

\paragraph{Exact range and comparison with the signature route.}
The theorem is valid for each specified product degeneration
meeting all the stated smoothness and transversality conditions.
It does not assert that an arbitrary hypersurface can be put
in that form by a permitted operation on its cycles.
For $n\ge5$ the conclusion here is codimension two, with
the already known complementary one-cycle degree; it does
not supply all middle codimensions.  A number-field variety
with this local coverage and the frozen coverage at all
other finite places has the full pairing on
$CH^2_{\mathrm{geom\,hom},\mathbf Q}\times
CH^{n-1}_{\mathrm{geom\,hom},\mathbf Q}$ by PH16.
No such global placewise hypothesis is omitted.

The new mechanism is the actual exceptional-divisor block
in \eqref{r18nf:triple-formula} and its elimination by
literal triple relations.  The broad signature criterion
S18 remains useful independently, but is unnecessary for
this theorem.  General $U_2$ and the residual homological
period map on arbitrary varieties remain open.

\paragraph{Dependencies and primary inputs.}
The resolution uses HP18's explicitly checked etale charts,
with the substitution justified above.  DG18 supplies the
unconditional degree-one lemma and its injection for $U_2$.
The separate cycle comparison uses CD17, CL17, classical
weak Lefschetz and Poincare duality in geometric $\ell$-adic
cohomology; it does not assume divisor Tate.  A primary
reference for the complete-intersection cohomology used here
is S.~Ghorpade and G.~Lachaud, \emph{Etale cohomology,
Lefschetz theorems and number of points of singular varieties
over finite fields}, corrected arXiv version 0808.2169,
Propositions 3.1--3.2; in the smooth case its nonmiddle
cohomology statement includes the required $H^0$ and $H^2$.
The remaining analytic
and height inputs are the final frozen ER17, VT17 and PH16
interfaces.  No historical priority claim is made.
